\section{The periodic Coulomb energy}\label{sec:setting}

\subsection{Definition}\label{ssec:energy}

A \emph{lattice} is a discrete subgroup $\Gamma\subset\R^3$ of rank three. We write $|\Gamma|$ for its covolume and
$\Gamma^*=\{q\in\R^3: q\cdot\gamma\in2\pi\Z\text{ for all }\gamma\in\Gamma\}$ for its dual lattice. The \emph{periodic Coulomb kernel}
of $\Gamma$ is the $\Gamma$-periodic, locally integrable function $V_\Gamma$ on $\R^3$, unique up to a null set, with Fourier
coefficients
\begin{equation}\label{eq:GF}
  \frac1{|\Gamma|}\int_{\R^3/\Gamma}V_\Gamma(x)\,e^{-iq\cdot x}\,dx=\begin{cases}4\pi/(|\Gamma|\,|q|^2),&q\in\Gamma^*\setminus\{0\},\\ 0,&q=0.\end{cases}
\end{equation}
It has a representative that is smooth on $\R^3\setminus\Gamma$ (Appendix~\ref{app:ewald}), which we use throughout; it is even,
and it is the solution of mean zero of
$-\Delta V_\Gamma=4\pi\bigl(\sum_{\gamma\in\Gamma}\delta_\gamma-|\Gamma|^{-1}\bigr)$: the potential of a unit point charge at each point of
$\Gamma$ in a uniform background of opposite total charge. The limit
\begin{equation}\label{eq:zeta}
  \zeta_\Gamma:=\lim_{x\to0}\Bigl(V_\Gamma(x)-\frac1{|x|}\Bigr)
\end{equation}
exists (Proposition~\ref{prop:ewald}); it is the potential that a charge feels from its own periodic images and the background.

A nonempty set $X\subset\R^3$ is a \emph{periodic configuration} if $X+\gamma=X$ for all $\gamma$ in some lattice $\Gamma$ and $X/\Gamma$ is
finite.
If $x_1,\dots,x_N$ are representatives of $X/\Gamma$, the \emph{energy per particle} of $X$ is
\begin{equation}\label{eq:edef}
  e(X):=\frac1N\Bigl(\sum_{1\le i<j\le N}V_\Gamma(x_i-x_j)+\frac N2\,\zeta_\Gamma\Bigr).
\end{equation}
This is the electrostatic energy per particle of unit point charges at the points of $X$ in a uniform neutralizing background
(the jellium, or one-component plasma, energy), with the infinite self-energies of the points removed; see
\cite[Section~IV]{Lewin2022} and \cite[Appendix~B]{RougerieSerfaty2016} for this definition and its relation to the thermodynamic
limit. For a lattice $\Lambda$, taking
$\Gamma=\Lambda$ gives $e(\Lambda)=\frac12\zeta_\Lambda$, the Madelung energy of $\Lambda$. The \emph{density} of $X$ is $N/|\Gamma|$.

\begin{proposition}[Ewald representation]\label{prop:ewald}
The limit \eqref{eq:zeta} exists, and for every $\alpha>0$ the quantity $Ne(X)$ of \eqref{eq:edef} equals
\begin{multline}\label{eq:ewald}
  \frac12\sum_{i,j=1}^N\ \sideset{}{'}\sum_{\gamma\in\Gamma}\frac{\erfc\bigl(\alpha|x_i-x_j+\gamma|\bigr)}{|x_i-x_j+\gamma|}
  +\frac{2\pi}{|\Gamma|}\sum_{q\in\Gamma^*\setminus\{0\}}\frac{e^{-|q|^2/4\alpha^2}}{|q|^2}\,\Bigl|\sum_{j=1}^Ne^{iq\cdot x_j}\Bigr|^2
  -\frac{\alpha N}{\sqrt\pi}-\frac{\pi N^2}{2\alpha^2|\Gamma|},
\end{multline}
where the prime means that $\gamma=0$ is omitted when $i=j$. Both series converge absolutely.
\end{proposition}

This is Ewald's formula \cite{Ewald1921}; a short proof from \eqref{eq:GF} is given in Appendix~\ref{app:ewald}. The next statement
collects the elementary invariances.

\begin{lemma}\label{lem:invariance}
\textup{(i)} $e(X)$ does not depend on the choice of the period lattice $\Gamma$ or of the representatives.
\textup{(ii)} $e(QX+v)=e(X)$ for every $Q\in\Or(3)$ and $v\in\R^3$.
\textup{(iii)} $e(tX)=t^{-1}e(X)$ for every $t>0$.
\end{lemma}

\begin{proof}
(ii) and (iii) follow from \eqref{eq:ewald}, which involves only the differences $x_i-x_j$ and is therefore invariant under
translations: for (ii) replace $\Gamma$ by $Q\Gamma$, and for (iii) replace $\Gamma$ by $t\Gamma$ and $\alpha$
by $\alpha/t$. For (i), independence of the representatives is clear from periodicity. The set of all periods of $X$ is a lattice
$P\supset\Gamma$, so two period lattices have a common sublattice of finite index, and it is enough to compare $\Gamma$ with a sublattice
$\Gamma'\subset\Gamma$ of index $k$. Evaluate \eqref{eq:ewald} for both with the same $\alpha$; the points for $\Gamma'$ are $x_j+\rho_l$,
where $\rho_1,\dots,\rho_k$ represent $\Gamma/\Gamma'$. For fixed $(i,l)$ and $j$, as $l'$ and $\gamma'\in\Gamma'$ vary, the vectors
$x_i+\rho_l-x_j-\rho_{l'}+\gamma'$ run once over $x_i-x_j+\Gamma$, the omitted term $\gamma=0$ for $i=j$ corresponding to
$(j,l')=(i,l)$, $\gamma'=0$; so the first sum is multiplied by $k$. For $q\in\Gamma'^*$ the structure factor is
$\sum_je^{iq\cdot x_j}\sum_le^{iq\cdot\rho_l}$, and $\sum_le^{iq\cdot\rho_l}$ equals $k$ if $q\in\Gamma^*$ and $0$ otherwise; since
$|\Gamma'|=k|\Gamma|$ the second term is multiplied by $k$. So are the last two terms, and $N$ is replaced by $kN$.
\end{proof}

By (iii) there is no loss in fixing the density. \emph{Throughout the paper all configurations have density $3/(4\pi)$}, so that the
volume per particle is that of the unit ball; energies are then expressed in the customary unit $e^2/a$, with $a$ the Wigner--Seitz
radius. At another density $\varrho$ all statements hold after the dilation by $(3/(4\pi\varrho))^{1/3}$.

\subsection{A series with nonnegative terms}\label{ssec:series}

For $a>0$ and $t>0$ let
\begin{equation}\label{eq:Gdef}
  \Ga(a,t):=\int_1^\infty e^{-tu}u^{a-1}\,du ,\qquad\text{so that}\qquad
  \Ga(\tfrac12,t)=\sqrt{\frac\pi t}\,\erfc\bigl(\sqrt t\bigr),\qquad \Ga(1,t)=\frac{e^{-t}}t .
\end{equation}
Each $\Ga(a,\cdot)$ is positive, decreasing and convex; indeed $\partial_t^k\Ga(a,t)=(-1)^k\Ga(a+k,t)$, and $\Ga(a,\cdot)$ is
logarithmically convex, being a superposition of exponentials. For a positive definite symmetric matrix $M$ and $v\in\R^3$ we write
$M[v]:=v^\T Mv$. For an integer $N\ge1$ put
\[
  \ell_N:=\Bigl(\frac{4\pi N}3\Bigr)^{1/3},\qquad\text{in particular}\qquad \lone=\Bigl(\frac{4\pi}3\Bigr)^{1/3},\quad \ltwo=\Bigl(\frac{8\pi}3\Bigr)^{1/3}=2.0309\ldots,
\]
the side of a cube that contains $N$ particles at density $3/(4\pi)$.

\begin{proposition}\label{prop:series}
Let $X$ be a periodic configuration of density $3/(4\pi)$ with period lattice $\Gamma=\ell_NR\,\Z^3$, $|\det R|=1$, and with $N$ points
$x_j=\ell_NRf_j$ per cell. Put $M=R^\T R$. Then
\begin{equation}\label{eq:seriesN}
  e(X)=\frac1{2N\ell_N}\biggl[\ \sum_{i,j=1}^N\ \sideset{}{'}\sum_{m\in\Z^3}\Ga\bigl(\tfrac12,\pi M[m+f_i-f_j]\bigr)
    +\sum_{m\in\Z^3\setminus\{0\}}\Bigl|\sum_{j=1}^Ne^{2\pi im\cdot f_j}\Bigr|^2\Ga\bigl(1,\pi M^{-1}[m]\bigr)-N^2-2N\biggr],
\end{equation}
where the prime omits $m=0$ when $i=j$. Every term of the two series is nonnegative.
\end{proposition}

\begin{proof}
Take $\alpha=\sqrt\pi/\ell_N$ in \eqref{eq:ewald}. A vector $x_i-x_j+\gamma$ has the form $\ell_NR(m+f_i-f_j)$, of length
$\ell_N t$ with $t^2=M[m+f_i-f_j]$, and $\erfc(\sqrt\pi t)/(\ell_Nt)=\ell_N^{-1}\Ga(\frac12,\pi t^2)$ by \eqref{eq:Gdef}. A dual vector is
$q=2\pi\ell_N^{-1}R^{-\T}m$, with $|q|^2=4\pi^2\ell_N^{-2}M^{-1}[m]$, so that
$\frac{2\pi}{|\Gamma|}|q|^{-2}e^{-|q|^2/4\alpha^2}=\frac1{2\ell_N}\Ga(1,\pi M^{-1}[m])$. Finally $\alpha N/\sqrt\pi=N/\ell_N$ and
$\pi N^2/(2\alpha^2|\Gamma|)=N^2/(2\ell_N)$.
\end{proof}

We use \eqref{eq:seriesN} for $N=1$ and $N=2$. For a lattice $\Lambda=\lone R\,\Z^3$,
\begin{equation}\label{eq:H1}
  e(\Lambda)=\frac{\Hone(M)-3}{2\lone},\qquad \Hone(M):=\sum_{m\ne0}\Bigl[\Ga\bigl(\tfrac12,\pi M[m]\bigr)+\Ga\bigl(1,\pi M^{-1}[m]\bigr)\Bigr].
\end{equation}
For $N=2$ we may take $f_1=0$ and $f_2=f\notin\Z^3$. The two cross terms $(i,j)=(1,2),(2,1)$ are equal, and the structure factor is
$|1+e^{2\pi im\cdot f}|^2=2+2\cos(2\pi m\cdot f)$. Hence
\begin{gather}\label{eq:H2}
  e(X)=\frac{\Htwo(M,f)-4}{2\ltwo},\\
  \Htwo(M,f):=\sum_{m\ne0}\Ga\bigl(\tfrac12,\pi M[m]\bigr)+\sum_{m\in\Z^3}\Ga\bigl(\tfrac12,\pi M[m+f]\bigr)
     +\sum_{m\ne0}\bigl(1+\cos2\pi m\cdot f\bigr)\,\Ga\bigl(1,\pi M^{-1}[m]\bigr).\notag
\end{gather}
We call the three sums of $\Htwo$ the \emph{same-coset}, the \emph{cross-coset} and the \emph{reciprocal} sum. Since all terms are
nonnegative, \emph{every partial sum of $\Htwo$ is a rigorous lower bound for it}; no estimate of an omitted tail is needed for lower
bounds on the energy. This elementary remark is used throughout Section~\ref{sec:cover}.

\begin{remark}
The quantity $\Hone(M)-3$ is the value at $s=\frac12$ of the analytic continuation of the Epstein zeta function
$\sum_{m\ne0}M[m]^{-s}$, written by Riemann's splitting of the Mellin integral of the theta function; see
\cite{SarnakStrombergsson2006,MuratovNovaga2026}. Formula \eqref{eq:H2} is the corresponding expression for a two-point basis.
\end{remark}

\begin{lemma}\label{lem:Hsym}
For $\gamma\in\GL(3,\Z)$, $n\in\Z^3$ and $f\notin\Z^3$,
\[
  \Htwo(\gamma^\T M\gamma,\gamma^{-1}f)=\Htwo(M,f),\qquad \Htwo(M,f+n)=\Htwo(M,f),\qquad \Htwo(M,-f)=\Htwo(M,f).
\]
\end{lemma}

\begin{proof}
Re-index $m\mapsto\gamma m$ in the first two sums and $m\mapsto\gamma^{-\T}m$ in the third; re-index $m\mapsto m+n$; re-index $m\mapsto-m$.
\end{proof}

\subsection{The class of configurations}\label{ssec:class}

\begin{definition}\label{def:class}
$\Cone$ is the set of lattices of density $3/(4\pi)$, and $\Ctwo$ is the set of all sets $X=(L+a)\cup(L+b)$ with $L\subset\R^3$ a lattice
of covolume $8\pi/3$ and $a,b\in\R^3$, $b-a\notin L$. We call the elements of $\Ctwo$ \emph{two-site configurations}, $L$ a \emph{period
lattice} and $L+a$, $L+b$ the two \emph{cosets}.
\end{definition}

Every lattice $\Lambda\in\Cone$ belongs to $\Ctwo$: if $L\subset\Lambda$ is any of its seven sublattices of index two and
$\tau\in\Lambda\setminus L$, then $\Lambda=L\cup(L+\tau)$. Conversely $X=L\cup(L+\tau)$ is a lattice exactly when $2\tau\in L$. Thus
$\Ctwo$ is the set of all periodic configurations of density $3/(4\pi)$ with at most two points per cell. It is invariant under
isometries.

\emph{Normalized coordinates.} After a translation, every $X\in\Ctwo$ is $\ltwo Y$ with
\begin{equation}\label{eq:Ynorm}
  Y=L\cup(L+\tau),\qquad L=R\,\Z^3,\quad \det R=1,\qquad \tau=Rf,\quad f\in\R^3\setminus\Z^3 ,
\end{equation}
and $e(X)=(\Htwo(R^\T R,f)-4)/(2\ltwo)$. We call $Y$ the \emph{normalized configuration}, $R$ a \emph{period matrix} and $f$ the
\emph{fractional offset}. The pair $(R,f)$ is determined by $X$ up to $R\mapsto QR\gamma$, $f\mapsto\pm\gamma^{-1}f+n$ with $Q\in\Or(3)$,
$\gamma\in\GL(3,\Z)$, $\det(Q\gamma)=1$, $n\in\Z^3$ (change of basis, of the representative of the second coset, and exchange of the two cosets) and, when
$X$ is a lattice, up to the choice of the sublattice of index two. Modulo isometries, $\Ctwo$ has eight parameters: five for the shape
of $L$ and three for $\tau$.

\subsection{The cubic lattices and their energies}\label{ssec:constants}

Let $\one=(1,1,1)$. The normalized body-centred and face-centred cubic lattices are
\begin{equation}\label{eq:B0F0}
  \Bz:=\Z^3\cup\bigl(\Z^3+\tfrac12\one\bigr),\qquad \Fz:=2^{-2/3}\bigl\{m\in\Z^3:\ m_1+m_2+m_3\in2\Z\bigr\} .
\end{equation}
Both have covolume $\frac12$, so $\ltwo\Bz$ and $\ltwo\Fz$ have density $3/(4\pi)$. We say that $X$ \emph{is a BCC lattice} (respectively
\emph{an FCC lattice}) if $X=Q(\ltwo\Bz)+v$ (respectively $X=Q(\ltwo\Fz)+v$) for some $Q\in\Or(3)$, $v\in\R^3$, and we put
\[
  \eB:=e(\ltwo\Bz),\qquad \eF:=e(\ltwo\Fz),\qquad \gap:=\eF-\eB .
\]

\begin{proposition}[Madelung energies]\label{prop:constants}
\begin{align*}
  \eB&=-0.89592925568176528916624546293829\ldots,\\
  \eF&=-0.89587361519510845989748779611028\ldots,\\
  \gap&=\phantom{-}5.5640486656829268757666828009\ldots\times10^{-5}.
\end{align*}
In particular $\gap>0$.
\end{proposition}

These values are classical \cite{Fuchs1935,ColdwellHorsfallMaradudin1960,Foldy1978}; the digits displayed here are certified, see
Section~\ref{ssec:constcert}. The relative difference of the two energies is $6.2\times10^{-5}$; the smallness of such differences is
the main difficulty of the problem.
